\section{Result and scope}
\label{sec:result}

Paper~16's five cascade schedules pass several of its finite necessary-condition
tests. That does not establish compatibility with earlier necessary
conditions. Here the relevant earlier result is Paper~2, Theorem~7.5
\cite{SecondPaper2026}: a contracting recursion with regular fixed-stage
motion cannot perform near-full-entropy transposition with sublinear
finite-control bits. The application requires the complete trace, not just
the stage that appears most complicated.

Throughout, $d=2^m$ is the matrix dimension, $b$ the entry width and
$L=d^2b$ the number of input bits. The two width regimes in our conclusion
are $b=1$ and $b=\max(1,\lceil\log_2d\rceil)=m$. All logarithms in
information quantities are base two. The five constructors are exactly
those of the published Rust crate at the revision recorded in
\cref{sec:verification}; their names are given in \cref{tab:bounds}.
For asymptotic statements, their loops and coordinate formulas are read
with exact integer arithmetic. The Rust executable realizes finite
representable instances; fixed-width integer overflow is not part of the
mathematical schedule family.

A \emph{completion} fixes the input-head moves, work-head moves, append
times and total duration to those of the constructor. Work tapes start
blank, all heads start at zero, and the output starts empty. At each tick,
an arbitrary deterministic local rule reads the old controller, raw input
symbol and all old scanned work symbols, then updates the controller,
writes the work symbols, and emits a bit when prescribed. Rules may depend
on $d$ and the public tick index. They may rewrite at \emph{every} visit,
including positioning steps marked \texttt{PassThrough} in the crate.
This is the unrestricted local-update interpretation of the schedule.
It includes completions that respect
the crate's more restrictive active-write annotations.

Fix a finite work alphabet $\Gamma$, of cardinality $A$, containing a
blank symbol. A completion has a nonempty finite controller set $Q$;
write $q=\lceil\log_2|Q|\rceil$. The initial controller value is public.
If the output equals raw row-major transpose on a nonempty set
$S\subseteq\{0,1\}^L$, write
\[
 \deficit=L-\log_2|S|.
\]
No assumption is made about correctness outside $S$, linearity of the
rules, independence of stored symbols, or uniform generation of the rules.

\begin{theorem}[C: exclusion of all five exact schedules]
\label{thm:cascades}
For every fixed finite work alphabet there exist $\delta>0$ and $d_0$
such that, for each of the five constructors, every dyadic $d\ge d_0$,
each of the two width regimes $b=1$ and $b=\max(1,\lceil\log_2d\rceil)$,
and every completion and nonempty correctness set $S$ as above,
\begin{equation}
\label{eq:main}
                 q+L-\log_2|S|\ge\delta L.
\end{equation}
The constants can be chosen simultaneously for all five constructors and
both widths. They are not evaluated numerically here.
\end{theorem}
\partformal{CascadeObstruction/C5Exclusion.lean}{CascadeObstruction.C5.uniform_c5_gap}{The
canonical-width C5 case is formalized for the declarative schedule in Lean.
Its Rust correspondence, the other four schedules, all bit-width cases and
the finite minimum across cases remain written proofs.}
The complete written proof is assembled in \cref{sec:bridges}, using the
source proofs in \cref{sec:schedules} and the cited Lean-verified obstruction
in \cref{sec:criterion}. For C5 at canonical width, the Lean theorem
discharges the compiler, clock and regularity obligations; its proof is
outlined in \cref{sec:c5-formal}. The verification boundary is recorded in
\cref{sec:verification}.

Full correctness gives $q\ge\delta L$. If $q=o(L)$, the ambient success
fraction is at most $2^{-(\delta-o(1))L}$. Thus fixed-state machines are
covered, but fixing the number of states is unnecessary. The number of
work tapes and alphabet are fixed; for these constructors there are two
or three work tapes. The complete schedule is oblivious even though its
data updates may be nonlinear and depend on previous data.

The statement is asymptotic along dyadic dimensions. Neither the smallest
legal constructor size nor a tested experimental size is an exclusion
threshold. The geometric estimates below hold for every positive integer
$b$, but the transposition conclusion is proved only in the two stated
width regimes. It does not decide unrestricted linear-time transposition,
question (T) of the series \cite{FirstPaper2026}.
